\chapter{The Zero-Free Quarter Disk and Complete Spectral Disjointness}
\label{chap:g018-zero-free-quarter-disk}

\status{SOH-G018 proved: $F$ has no zeros in $|w|\le1/4$, hence the canonical negative inversion maps no $F$ root to an $F$ root and no xi zero to an xi zero.}

\section{From finite exceptions to no exceptions}

Chapters~\ref{chap:quotient-negative-inversion-zero-set-no-go}--\ref{chap:g017-fixed-point-exclusion} progressively restricted the possible overlap between the root set of the even quotient $F$ and its image under
\[
 J(w)=\frac{1}{16w}.
\]
The present chapter closes that frontier.  The key is a zero-free disk whose boundary radius is exactly the self-reciprocal radius of $J$.

Recall
\[
 \xi\!\left(\frac12+z\right)=F(z^2),
 \qquad
 F(w)=\sum_{n\ge0}a_nw^n,
\]
with
\[
 a_n>0\qquad(n\ge0)
\]
from the positive Riemann-kernel coefficient theorem.

\section{The exact quarter-point value}

At $w=1/4$ choose $z=1/2$.  Then
\[
 F\!\left(\frac14\right)=\xi(1).
\]
For
\[
 \xi(s)=\frac12s(s-1)\pi^{-s/2}\Gamma(s/2)\zeta(s),
\]
the simple pole of $\zeta(s)$ at $s=1$ has residue one, hence
\begin{equation}
 \boxed{F(1/4)=\xi(1)=\frac12}.
 \label{eq:g018-quarter-value}
\end{equation}

\section{An elementary lower bound for $F(0)$}

At the centre,
\[
 F(0)=\xi\!\left(\frac12\right)
 =-\frac18\pi^{-1/4}\Gamma\!\left(\frac14\right)\zeta\!\left(\frac12\right).
\]
The Dirichlet eta identity gives
\[
 \eta(s)=\bigl(1-2^{1-s}\bigr)\zeta(s),
\]
so at $s=1/2$,
\[
 -\zeta\!\left(\frac12\right)
 =(1+\sqrt2)\eta\!\left(\frac12\right).
\]
The alternating series for $\eta(1/2)$ has decreasing terms.  Its fourth partial sum is therefore a strict lower bound:
\[
 \eta\!\left(\frac12\right)
 >
 1-\frac1{\sqrt2}+\frac1{\sqrt3}-\frac12
 =
 \frac12-\frac1{\sqrt2}+\frac1{\sqrt3}.
\]

For the gamma factor, $e^{-t}>1-t$ for $0<t\le1$, and therefore
\[
 \Gamma\!\left(\frac14\right)
 >
 \int_0^1t^{-3/4}(1-t)\,dt
 =\frac{16}{5}.
\]
Also $\pi<4$ gives
\[
 \pi^{-1/4}>\frac1{\sqrt2}.
\]
Consequently
\[
 F(0)>L,
\]
where
\[
 L:=\frac{2}{5\sqrt2}(1+\sqrt2)
 \left(\frac12-\frac1{\sqrt2}+\frac1{\sqrt3}\right).
\]

The comparison with $1/4$ can be made exact.  Simplification gives
\[
 60\left(L-\frac14\right)
 =4\sqrt6+8\sqrt3-15-6\sqrt2.
\]
Both sides of
\[
 4\sqrt6+8\sqrt3>15+6\sqrt2
\]
are positive.  Their squared difference equals
\[
 -9+12\sqrt2.
\]
Since $\sqrt2>3/4$, equivalently $32>9$, this difference is positive.  Hence
\begin{equation}
 \boxed{F(0)>L>\frac14}.
 \label{eq:g018-centre-bound}
\end{equation}

\section{Zero-free disk theorem}

Let $a_0=F(0)$.  For $|w|\le1/4$, positivity of all Taylor coefficients and absolute convergence give
\[
 \begin{aligned}
 |F(w)-a_0|
 &\le\sum_{n\ge1}a_n|w|^n\\
 &\le\sum_{n\ge1}a_n4^{-n}\\
 &=F(1/4)-F(0)\\
 &=\frac12-a_0.
 \end{aligned}
\]
By \eqref{eq:g018-centre-bound}, $a_0>1/4$, so $1/2-a_0<a_0$.  Therefore
\[
 |F(w)|
 \ge a_0-|F(w)-a_0|
 \ge2a_0-\frac12
 >2L-\frac12
 >0.
\]

\begin{theorem}[SOH-G018 zero-free quarter disk]
\label{thm:g018-zero-free-quarter-disk}
The even quotient $F$ has no zeros in the closed disk
\[
 \boxed{|w|\le\frac14}.
\]
Equivalently, every root of $F$ satisfies $|w|>1/4$.
\end{theorem}

\section{Complete negative-inversion exclusion}

Suppose that $w$ and $J(w)$ were both roots of $F$.  Theorem~\ref{thm:g018-zero-free-quarter-disk} gives
\[
 |w|>\frac14.
\]
But then
\[
 |J(w)|=\frac{1}{16|w|}<\frac14,
\]
so the same theorem gives $F(J(w))\ne0$, a contradiction.

Thus the paired quotient set of Chapter~\ref{chap:quotient-negative-inversion-zero-set-no-go} is not merely finite:
\begin{equation}
 \boxed{P_J=\varnothing}.
 \label{eq:g018-pj-empty}
\end{equation}
Equivalently,
\[
 \boxed{Z_F\cap J(Z_F)=\varnothing}.
\]

Chapter~\ref{chap:paired-spectrum-quotient-correspondence} proves the exact relation
\[
 q(P_N)=P_J,
 \qquad
 q(s)=\left(s-\frac12\right)^2.
\]
Equation~\eqref{eq:g018-pj-empty} therefore implies
\begin{equation}
 \boxed{P_N=\varnothing}.
 \label{eq:g018-pn-empty}
\end{equation}
In explicit xi language,
\[
 \boxed{
 \xi(\rho)=0
 \quad\Longrightarrow\quad
 \xi\!\left(\frac{\rho-1}{2\rho-1}\right)\ne0
 }
\]
for every xi zero $\rho$.

\section{What has been closed}

G014 proved finiteness of the possible paired xi-zero set.  G015 proved finiteness of the quotient paired-root set.  G016 identified those finite sets through an exact two-to-one quotient.  G017 removed all fixed exceptional roots.  G018 proves that no exceptional two-cycle exists either:
\[
 \boxed{P_J=P_N=\varnothing}.
\]
The canonical negative inversion is therefore spectrally disjoint from the xi zero set.

\section{Boundary of the theorem}

The result is a global statement about one exact projective involution and the xi/$F$ zero sets.  It does not prove real-rootedness of $F$, PF$\infty$, SOH-G003, or the Riemann Hypothesis.  In particular, spectral disjointness under $N_s$ is logically different from the critical-line statement required by RH.
